% Section 7: Amplification Loop
\label{sec:amplification}

\subsection{Joint System Definition}

The three preceding sections each demonstrate that AI signal correlation creates fragility in isolation. Each also shows why its channel alone provides an incomplete assessment. We now formalise why. Channel 3 determines the effective number of independent liquidity providers, which amplifies effective signal correlation in Channel 1. Channel 1 determines the crisis threshold, which reduces the option value of private information in Channel 2. Channel 2 determines the fraction of privately informed agents, which feeds back into Channel 3 by altering market-maker reliance on the common AI signal. The system is circular. This section formalises the joint dynamics as a fixed-point problem.

The state vector of the amplification loop is $v = (\rho_{\text{eff}}, \theta^*, N_{\text{eff}}) \in K$, where the domain is
\begin{equation}\label{eq:domain}
  K = [\rho, 1] \times [\theta_L, \theta_H] \times [1, N].
\end{equation}
Here $\rho$ is the exogenous AI signal correlation (a parameter of the system), $[\theta_L, \theta_H]$ are the Goldstein-Pauzner (2005) dominance region boundaries, and $N \geq 2$ is the number of market makers. The set $K$ is a closed bounded rectangle in $\mathbb{R}^3$, hence compact and convex.

The effective correlation $\rho_{\text{eff}} \in [\rho, 1]$ incorporates both exogenous AI model correlation and endogenous amplification from correlated liquidity withdrawal. The crisis threshold $\theta^* \in [\theta_L, \theta_H]$ is determined by the Goldstein-Pauzner indifference condition with $\rho_{\text{eff}}$ as input. The effective market-maker count $N_{\text{eff}} \in [1, N]$ depends on $\rho$, $N$, and the fraction of privately informed agents $\mu_I$ from Channel 2.

The three equilibrium mappings connect these state variables.

\prooflabel{Mapping 1 (Channel 3 to Channel 1): $g_1$.} When the effective number of independent liquidity providers is low, correlated market-maker withdrawal amplifies price dislocations, which function as an additional common signal that raises effective correlation in the coordination game.

\medskip\noindent\textbf{Assumption A2} (Aggregation Form).\label{ass:g1}
\emph{The effective signal correlation in the coordination game is}
\begin{equation}\label{eq:g1}
  \rho_{\text{eff}} = g_1(\rho, N_{\text{eff}}) = 1 - (1-\rho)\frac{N_{\text{eff}}}{N}.
\end{equation}

\noindent The formula in~\eqref{eq:g1} is a convex combination between $\rho$ (the exogenous level when $N_{\text{eff}} = N$) and $1$ (the limit when $N_{\text{eff}} = 1$). Each unit reduction in $N_{\text{eff}}$ adds $(1-\rho)/N$ to effective correlation, reflecting one fewer independent source of liquidity. Any alternative aggregation form satisfying the same boundary conditions, monotonicity, and continuity yields qualitatively identical results; the linear form is chosen for tractability.

\prooflabel{Mapping 2 (Channel 1 to Channel 2): $g_2$.} The crisis threshold $\theta^*$ determines the crisis probability and hence the option value of private fundamental research. The equilibrium fraction of privately informed agents $\mu_I^*$ solves the modified Grossman-Stiglitz indifference condition~\eqref{eq:gs-indifference}:
\begin{equation}\label{eq:g2}
  \mu_I^* = g_2(\theta^*, \rho_{\text{eff}}, c_P): \quad \pi_P(\mu_I^*) - \pi_A(\rho_{\text{eff}}, 1-\mu_I^*) = \frac{c_P}{1-\theta^*}.
\end{equation}
A higher crisis threshold $\theta^*$ raises the effective cost $c_P/(1-\theta^*)$ in~\eqref{eq:g2}, reducing $\mu_I^*$: more likely crises destroy the option value of private information.

\prooflabel{Mapping 3 (Channel 2 to Channel 3): $g_3$.} As $\mu_I^*$ falls, market makers rely more heavily on the common AI signal for pricing and inventory management, reducing the effective number of independent liquidity providers:
\begin{equation}\label{eq:g3}
  N_{\text{eff}} = g_3(\mu_I, \rho, N) = \frac{N}{1+(N-1)\rho(1-\mu_I)^2},
\end{equation}
where $(1-\mu_I)^2$ in~\eqref{eq:g3} captures the fraction of market-maker signal variance that is correlated. At $\mu_I = 0$ (all AI): $g_3$ reduces to the baseline~\eqref{eq:neff}. At $\mu_I = 1$ (all private): $N_{\text{eff}} = N$ (full independence).

\subsection{Fixed-Point Existence}

The composite operator $T: K \to K$ maps $v = (\rho_{\text{eff}}, \theta^*, N_{\text{eff}})$ through the four-step composition:
\begin{equation}\label{eq:composite-T}
\begin{aligned}
  T_3(v) &= g_3\!\left(g_2(\theta^*, \rho_{\text{eff}}, c_P),\; \rho,\; N\right), \\
  T_1(v) &= g_1(\rho, T_3(v)) = 1 - (1-\rho)\frac{T_3(v)}{N}, \\
  T_2(v) &= \theta^*(T_1(v)),
\end{aligned}
\end{equation}
where $\theta^*(\cdot)$ denotes the Goldstein-Pauzner crisis threshold function of Channel 1. The composition in~\eqref{eq:composite-T} makes the circular structure explicit: Channel 3 feeds Channel 1, which feeds Channel 2, which feeds back to Channel 3.

\begin{proposition}[Fixed-Point Existence]\label{prop:fixed-point-existence}
  The composite operator $T: K \to K$ is continuous and maps the compact convex set $K$ defined in~\eqref{eq:domain} into itself. By Brouwer's fixed-point theorem, $T$ admits at least one fixed point $(\rho_{\text{eff}}^*, \theta^{**}, N_{\text{eff}}^*) \in K$.

  \prooflabel{Proof.} $K$ is compact and convex. $T$ maps $K$ into $K$: $T_3(v) \in [1, N]$ because $\mu_I \in [0,1]$; $T_1(v) \in [\rho, 1]$ because $T_3(v) \in [1, N]$; $T_2(v) \in [\theta_L, \theta_H]$ by the dominance argument. Each component ($g_1$ linear, $g_3$ rational, $g_2$ and $\theta^*$ $C^1$ by the implicit function theorem) is continuous, so $T$ is continuous. Brouwer's theorem applies. $\square$
\end{proposition}

\subsection{Jacobian Analysis}

At an interior fixed point $v^* = (\rho_{\text{eff}}^*, \theta^{**}, N_{\text{eff}}^*)$ with $\rho_{\text{eff}}^* < \rho_1^*$, five partial derivatives govern the strength of each link in the feedback loop: the Channel 2 responses $a \equiv d\mu_I^*/d\theta^* < 0$ (see~\eqref{eq:coeff-a}) and $b \equiv d\mu_I^*/d\rho_{\text{eff}} > 0$ (see~\eqref{eq:coeff-b}), the Channel 3 liquidity sensitivity $m \equiv dg_3/d\mu_I > 0$ (see~\eqref{eq:coeff-m}), the Channel 1 crisis sensitivity $h \equiv d\theta^*/d\rho_{\text{eff}} > 0$ (see~\eqref{eq:coeff-h}), and the $g_1$ amplification coefficient $w \equiv (1-\rho)/N > 0$ (see~\eqref{eq:coeff-w}). The full derivation of these coefficients is in Appendix~A. The Jacobian matrix of $T$ at $v^*$ is
\begin{equation}\label{eq:jacobian}
  DT = \begin{pmatrix} -wmb & -wm|a| & 0 \\ -hwmb & hwm|a| & 0 \\ mb & -m|a| & 0 \end{pmatrix},
\end{equation}
where the third column is zero because the updated state vector $(T_1(v), T_2(v), T_3(v))$ depends on the current $(\rho_{\text{eff}}, \theta^*)$ but not on the current $N_{\text{eff}}$ (which enters only through $g_1$, with $N_{\text{eff}}'$ computed from the current $\theta^*$ and $\rho_{\text{eff}}$). The structure of~\eqref{eq:jacobian} reveals that the system has rank two: one eigenvalue is zero, and stability depends entirely on the non-trivial eigenvalue
\begin{equation}\label{eq:eigenvalue}
  \ell_1 = w \cdot m \cdot (h|a| - b).
\end{equation}
The sign of $\ell_1$ in~\eqref{eq:eigenvalue} depends on whether the cross-channel destabilising effect $h|a|$ (the product of Channel 1's sensitivity to $\rho_{\text{eff}}$ and Channel 2's sensitivity to $\theta^*$) exceeds the within-channel stabilising effect $b$ (the substitution from AI to private information as AI rents collapse).

The positive feedback loop is concrete. A rise in $\rho_{\text{eff}}$ raises $\theta^*$ through Channel 1 ($h > 0$). Higher $\theta^*$ lowers $\mu_I^*$ through Channel 2 ($a < 0$). Lower $\mu_I^*$ lowers $N_{\text{eff}}$ through Channel 3 ($m > 0$). Lower $N_{\text{eff}}$ raises $\rho_{\text{eff}}$ through $g_1$ ($w > 0$). Every step preserves the sign of the initial perturbation. The loop closes. The interior fixed point is stable when $|\ell_1| < 1$ and unstable when $|\ell_1| > 1$. The bifurcation occurs at $|\ell_1| = 1$.

\subsection{Bifurcation and the Safety Illusion}

\begin{proposition}[Amplification Bifurcation]\label{prop:bifurcation}
  Let $\rho_1^* = 1/(1+\alpha_{\text{SC}}^2)$ denote the Channel 1 uniqueness boundary (Proposition~\ref{prop:uniqueness-boundary}), $\rho_2^*$ the Channel 2 threshold at which revelatory price efficiency begins declining (Proposition~\ref{prop:rpe}), and $\rho_3^* = \rho^{**}$ the Channel 3 no-equilibrium threshold (Proposition~\ref{prop:rho-star-star}). Define $\rho^*$ as the value of the exogenous correlation $\rho$ at which the spectral radius of $DT$ at the interior fixed point crosses unity. Then
  \begin{equation}\label{eq:bifurcation}
    \rho^* < \min(\rho_1^*, \rho_2^*, \rho_3^*).
  \end{equation}

  \prooflabel{Proof sketch.} We proceed in three steps.

  \prooflabel{Step 1 (Individual channel stability).} For $\rho < \rho_i^*$ considered individually, each channel maintains a stable interior equilibrium. With cross-channel feedback severed (each channel's output feeds forward but does not return), the Jacobian is zero and the spectral radius is zero.

  \emph{Step 2 (Divergence of $h$ and positive spectral radius).} In the coupled system, the non-trivial eigenvalue is $\ell_1 = wm(h|a|-b)$. The Goldstein-Pauzner equilibrium system undergoes a saddle-node bifurcation at $\rho_{\text{eff}} = \rho_1^*$, where the $2 \times 2$ Jacobian of the equilibrium conditions satisfies $\det(J) = 0$ (see~\eqref{eq:h-divergence} in Appendix~B). Verifying the three non-degeneracy conditions of the saddle-node normal form (Guckenheimer and Holmes, 1983, Theorem~3.4.1), the local expansion~\eqref{eq:saddle-expansion} yields $h = d\theta^*/d\rho_{\text{eff}} = C/(2\sqrt{\rho_1^* - \rho_{\text{eff}}}) + O(1) \to \infty$ as $\rho_{\text{eff}} \to \rho_1^{*-}$. Since $|a|$ is bounded away from zero and $b$ is bounded, the divergence of $h$ ensures $h|a| > b$ for $\rho_{\text{eff}}$ sufficiently close to $\rho_1^*$, confirming $\ell_1 > 0$.

  \prooflabel{Step 3 (Strict inequality $\rho^* < \min_i \rho_i^*$).} Amplification from $g_1$ implies $\rho_{\text{eff}}^* > \rho$ whenever $N_{\text{eff}}^* < N$. Define $\rho^*$ as the solution to $\rho_{\text{eff}}^*(\rho^*) = \rho_1^*$; since $\rho_{\text{eff}}^* > \rho$, we obtain $\rho^* < \rho_1^*$. For Channels 2 and 3: the coupled system's effective correlation $\rho_{\text{eff}}^* > \rho$ and elevated crisis threshold $\theta^{**} > \theta^*(\rho)$ both tighten the indifference and participation constraints beyond what exogenous $\rho$ alone produces, ensuring $\rho^* < \rho_2^*$ and $\rho^* < \rho_3^*$. $\square$
\end{proposition}

The inequality~\eqref{eq:bifurcation} is the paper's centrepiece. The integrated three-channel system becomes fragile at a strictly lower level of AI signal correlation than any individual channel predicts. The result is not merely quantitative. It is structural: there exists a non-empty interval $(\rho^*, \min_i \rho_i^*)$ in which every individual channel assessment is favourable but the integrated system has already crossed its bifurcation threshold. What follows is the most important implication of the analysis.

\begin{proposition}[Local Uniqueness Below the Threshold]\label{prop:local-uniqueness}
  For $\rho < \rho^*$, the spectral radius of $DT$ at the interior fixed point satisfies $|\ell_1| < 1$. Since $\ell_1$ is the only non-zero eigenvalue, the spectral radius $\varrho(DT) = |\ell_1| < 1$. By the Ostrowski theorem (see, e.g., Ortega and Rheinboldt, 1970, Theorem 10.1.3), there exists a matrix norm $\|\cdot\|$ and a neighbourhood $U \subset K$ of the fixed point such that $\|DT(v)\| < 1$ for all $v \in U$, making $T$ a local contraction. By the Banach fixed-point theorem, the interior fixed point is locally unique within $U$.

  \prooflabel{Remark.} Global uniqueness is not established. The Brouwer theorem guarantees existence but not uniqueness globally. For $\rho > \rho^*$, the stable interior fixed point may be absorbed by the unstable manifold, leaving only a corner fixed point at $(1-(1-\rho)/N,\; \theta_H,\; 1)$, corresponding to maximal effective correlation, worst-case crisis threshold, and a single effective market maker.
\end{proposition}

\begin{corollary}[Safety Illusion]\label{cor:safety-illusion}
  Let $\rho \in (\rho^*, \min(\rho_1^*, \rho_2^*, \rho_3^*))$. Then:

  \emph{(i)} A regulator examining Channel 1 alone observes $\rho < \rho_1^*$ and concludes: the coordination game has a unique threshold equilibrium; no self-fulfilling crises are possible.

  \emph{(ii)} A regulator examining Channel 2 alone observes $\rho < \rho_2^*$ and concludes: the information acquisition equilibrium is interior and stable; price informativeness is adequate.

  \emph{(iii)} A regulator examining Channel 3 alone observes $\rho < \rho_3^*$ and concludes: the market-making equilibrium exists with finite spreads; market liquidity is sufficient.

  \emph{(iv)} The integrated system is fragile: $\rho > \rho^*$ implies the spectral radius of $DT$ exceeds unity, the stable interior fixed point is lost, and the positive feedback loop amplifies perturbations.

  \prooflabel{Proof.} Parts (i)--(iii) follow from the definitions of the individual thresholds $\rho_i^*$. Part (iv) follows from Proposition~\ref{prop:bifurcation}: $\rho > \rho^*$ implies $|\ell_1| > 1$. $\square$
\end{corollary}

The safety illusion is not an error of judgement. It is a structural property of the system. Each channel in isolation is genuinely stable at $\rho \in (\rho^*, \min_i \rho_i^*)$. The fragility arises from cross-channel positive feedback that no single-channel regulator can observe. \citet{danielsson2022} warned qualitatively that AI model monoculture creates systemic risk through multiple channels. Corollary~\ref{cor:safety-illusion} makes the warning precise: the interaction of the channels, not any single channel, is the source of the most dangerous systemic risk.

The width of the safety illusion interval $\min_i \rho_i^* - \rho^*$ increases with stronger strategic complementarity $\alpha_{\text{SC}}$, more market makers $N$, and higher cost of private information $c_P$. In the limiting case $N \to \infty$ with fixed $\alpha_{\text{SC}} > 0$, the bifurcation threshold $\rho^* \to 0$: even infinitesimal correlation is amplified through infinitely many feedback paths. With many correlated market makers, the usual diversification benefit of having multiple liquidity providers is entirely lost. The amplification loop can activate at arbitrarily low levels of AI model similarity.

The amplification result has a direct regulatory implication. Stress tests that evaluate coordination risk, information quality, and market liquidity separately will systematically underestimate systemic fragility when AI model homogeneity is the common driver across all three channels. The question is whether the market can avoid the safety illusion on its own, or whether regulatory intervention is necessary. Section~\ref{sec:extensions} answers this question by showing that competitive AI adoption endogenously drives $\rho$ above $\rho^*$, creating a prisoner's dilemma that only a diversity mandate can resolve.
% --- Clarity pass complete: 2026-03-10 ---
% --- Flow pass complete: 2026-03-11 ---
% --- Technical audit complete: 2026-03-11 ---
